\chapter{पूर्णांक अंकगणित}\label{ch:b1:arith}

अंकगणित --- $\Z$ में \omterm{def:b1:arith:divides}{विभाज्यता} का अध्ययन --- उच्चतर माध्यमिक खंड में आरंभ
हुआ था। यह अध्याय उसे यूक्लिडीय भाग से पूरी उपपत्तियों के साथ नए सिरे से
खड़ा करता है: \omterm{thm:b1:arith:gcd}{महत्तम समापवर्तक} और \omterm{met:b1:arith:euclid}{यूक्लिडीय कलनविधि}, बेज़ू सर्वसमिका और गाउस
की प्रमेयिका, \omterm{thm:b1:arith:fta}{अभाज्य गुणनखंडन}, तथा फर्मा की लघु प्रमेय तक सर्वांगसमताओं का
कलन। अपने आकर्षण से आगे, यह सामग्री वह आदर्श है जिसका अनुकरण
\cref{ch:b1:poly} बहुपदों के लिए करता है।

\section{विभाज्यता और यूक्लिडीय भाग}

\begin{definition}[विभाज्यता]\label{def:b1:arith:divides}
$a, b \in \Z$ के लिए, $b$ $a$ को \emph{विभाजित करता है} (लिखा जाता है $b
\mid a$)\index{विभाज्यता} जब किसी $q \in \Z$ के लिए $a = bq$ हो। मूल परिणाम:
यदि $b \mid a$ और $b \mid a'$, तो सभी $u, v \in \Z$ के लिए $b \mid (ua +
va')$; यदि $b \mid a$ और $a \neq 0$, तो $\abs b \leq \abs a$; और $a \mid b$
तथा $b \mid a$ साथ मिलकर $b = \pm a$ पर बाध्य कर देते हैं।
\end{definition}

\begin{theorem}[यूक्लिडीय भाग]\label{thm:b1:arith:division}
सभी $a \in \Z$ और $b \in \N^*$ के लिए ठीक एक युग्म $(q, r) \in \Z \times \N$
ऐसा है कि
\[
a = bq + r, \qquad 0 \leq r < b .
\]
\end{theorem}

\begin{proof}
\emph{अस्तित्व।} \omterm{def:b1:logic:sets}{समुच्चय} $A = \{a - bk : k \in \Z\} \cap \N$ $\N$ का अरिक्त
उपसमुच्चय है ($k = -\abs a$ लीजिए: $a + b\abs a \geq a + \abs a \geq 0$)।
मान लीजिए $r = a - bq$ उसका न्यूनतम अवयव है। यदि $r \geq b$, तो $r - b = a -
b(q+1)$ $A$ का उससे छोटा अवयव होता: विरोधाभास। अतः $0 \leq r < b$।

\emph{अद्वितीयता।} यदि $0 \leq r, r' < b$ के साथ $bq + r = bq' + r'$, तो
$b(q - q') = r' - r$ और $\abs{r' - r} < b$: बाएँ पक्ष में $b$ का गुणज $0$
होना ही चाहिए, अतः $q = q'$ और $r = r'$।
\end{proof}

\begin{example}[बार-बार भाग देकर स्थानीय मान पद्धति]\label{ex:b1:arith:base7}
$2026$ को आधार $7$ में लिखिए। $7$ से बार-बार भाग दीजिए और शेषफल रखते जाइए:
\[
2026 = 7 \times 289 + 3, \quad
289 = 7 \times 41 + 2, \quad
41 = 7 \times 5 + 6, \quad
5 = 7 \times 0 + 5 .
\]
शेषफलों को अंतिम से पहले तक पढ़िए: $2026 = (5\,6\,2\,3)_7$। जाँच: $5 \times
343 + 6 \times 49 + 2 \times 7 + 3 = 1715 + 294 + 14 + 3 = 2026$। यूक्लिडीय
भाग की अद्वितीयता ही हर अंक को \emph{अनिवार्य} बनाती है: हर पद पर शेषफल
$\intint06$ का एकमात्र ऐसा पूर्णांक है जो वर्तमान मान के $7$ के सापेक्ष
सर्वांगसम हो, अतः आधार-$7$ लेखन अद्वितीय है --- यही तथ्य चुपचाप तब प्रयुक्त
होता है जब सप्ताहांत समस्या ``आधार $p$ में $n$ के अंकों'' को साधती है।
\end{example}

\section{महत्तम समापवर्तक}

\begin{theorem}[$\Z$ के उपसमूह; महत्तम समापवर्तक का अस्तित्व]\label{thm:b1:arith:gcd}
\begin{enumerate}
  \item $(\Z, +)$ का प्रत्येक उपसमूह किसी अद्वितीय $n \in \N$ के लिए $n\Z =
  \{nk : k \in \Z\}$ रूप का है।
  \item जो $a, b \in \Z$ दोनों शून्य न हों, उनके लिए \omterm{def:b1:logic:sets}{समुच्चय} $a\Z + b\Z =
  \{au + bv : u, v \in \Z\}$ $\Z$ का उपसमूह है, अतः किसी अद्वितीय $d \in
  \N^*$ के लिए $d\,\Z$ के बराबर है। यही $d$ \emph{महत्तम
  समापवर्तक}\index{महत्तम समापवर्तक} $\gcd(a, b)$ है: वह $a$ और $b$ को
  \omterm{def:b1:arith:divides}{विभाजित करता है}, और $a$ तथा $b$ का प्रत्येक उभयनिष्ठ भाजक $d$ को \omterm{def:b1:arith:divides}{विभाजित
  करता है}।
\end{enumerate}
\end{theorem}

\begin{proof}
(1) मान लीजिए $H \subseteq \Z$ उपसमूह है (अरिक्त, घटाव के अंतर्गत स्थायी;
औपचारिक परिभाषा \cref{ch:b1:structures} में है, और केवल यही दो गुणधर्म
प्रयुक्त होते हैं)। यदि $H = \{0\}$, तो $n = 0$ लीजिए। अन्यथा $H$ में कोई
अशून्य अवयव और उसका विपरीत भी है, अतः उसमें एक लघुतम पूर्णतः धनात्मक अवयव
$n$ है। तब $n\Z \subseteq H$। $x \in H$ के लिए $0 \leq r < n$
(\cref{thm:b1:arith:division}) के साथ $x = nq + r$ लिखिए; $r = x - nq \in
H$, और $n$ की लघुतमता $r = 0$ पर बाध्य कर देती है: $x \in n\Z$। अद्वितीयता:
$n$ $n\Z$ का न्यूनतम धनात्मक अवयव है।

(2) $a\Z + b\Z$ में $0$ है और वह घटाव के अंतर्गत स्थायी है, अतः वह $d \geq
1$ के साथ $d\Z$ है (उसमें $a$ या $b$ अशून्य है)। चूँकि $a, b \in d\Z$, $d$
दोनों को \omterm{def:b1:arith:divides}{विभाजित करता है}। और यदि $c$ $a$ तथा $b$ को \omterm{def:b1:arith:divides}{विभाजित करता है}, तो $c$
प्रत्येक $au + bv$ को \omterm{def:b1:arith:divides}{विभाजित करता है} --- विशेष रूप से $c \mid d$ को,
क्योंकि $d \in a\Z + b\Z$। यही घोषित गुणधर्म है (और इससे $\abs c \leq d$
निकलता है, अतः $d$ \emph{महत्तम} समापवर्तक कहलाने का अधिकारी है)।
\end{proof}

\begin{corollary}[बेज़ू सर्वसमिका]\label{cor:b1:arith:bezout}
जो $a, b$ दोनों शून्य न हों, उनके लिए ऐसे $u, v \in \Z$ विद्यमान हैं कि
\[
au + bv = \gcd(a, b) .
\]
विशेष रूप से ($\gcd(a,b) = 1$, अर्थात् \emph{सहअभाज्य}\index{सहअभाज्य}
स्थिति): $a$ और $b$ \omterm{cor:b1:arith:bezout}{सहअभाज्य} हैं यदि और केवल यदि $au + bv = 1$ का कोई हल हो।
\end{corollary}

\begin{proof}
$\gcd(a,b) = d \in d\Z = a\Z + b\Z$। तुल्यता के लिए: यदि $\gcd(a,b) = 1$, तो
बेज़ू हल दे देता है; विलोमतः $au + bv = 1$ $a, b$ के प्रत्येक उभयनिष्ठ भाजक
को $1$ विभाजित करने पर बाध्य कर देता है।
\end{proof}

\begin{method}[यूक्लिडीय कलनविधि, विस्तारित]\label{met:b1:arith:euclid}
$\gcd(a, b)$ ($a > b > 0$) संगणित करने के लिए: $a = bq + r$ का भाग दीजिए;
फिर $\gcd(a, b) = \gcd(b, r)$ ($(a,b)$ और $(b,r)$ के उभयनिष्ठ भाजक एक ही
हैं, क्योंकि $r = a - bq$); शेषफल $0$ होने तक दोहराइए; अंतिम अशून्य शेषफल ही
\omterm{thm:b1:arith:gcd}{महत्तम समापवर्तक} है\index{यूक्लिडीय कलनविधि}। भागों को उलटी दिशा में चलाने
से (या नीचे उतरते समय गुणांक बनाए रखने से) एक बेज़ू युग्म $(u, v)$ मिल जाता
है।
\end{method}

\begin{example}\label{ex:b1:arith:euclid}
$\gcd(120, 23)$: $120 = 5 \times 23 + 5$; $23 = 4 \times 5 + 3$; $5 =
1\times 3 + 2$; $3 = 1 \times 2 + 1$; $2 = 2 \times 1 + 0$। अतः $\gcd = 1$।
उलटी दिशा में:
\begin{align*}
1 &= 3 - 2 = 3 - (5 - 3) = 2\times 3 - 5 = 2(23 - 4\times 5) - 5 \\
&= 2 \times 23 - 9 \times 5 = 2\times 23 - 9(120 - 5\times 23)
= 47 \times 23 - 9 \times 120 .
\end{align*}
जाँच: $47 \times 23 = 1081$, $9 \times 120 = 1080$।
\end{example}

\begin{theorem}[गाउस की प्रमेयिका और उसके परिणाम]\label{thm:b1:arith:gauss}
मान लीजिए $a, b, c \in \Z$।
\begin{enumerate}
  \item (गाउस की प्रमेयिका\index{गाउस की प्रमेयिका}) यदि $a \mid bc$ और
  $\gcd(a, b) = 1$, तो $a \mid c$।
  \item यदि $a \mid c$, $b \mid c$ और $\gcd(a,b) = 1$, तो $ab \mid c$।
  \item यदि $\gcd(a, b) = \gcd(a, c) = 1$, तो $\gcd(a, bc) = 1$।
\end{enumerate}
\end{theorem}

\begin{proof}
(1) बेज़ू: $au + bv = 1$। $c$ से गुणा कीजिए: $acu + bcv = c$। दोनों पद $a$
से विभाज्य हैं (दूसरा इसलिए कि $a \mid bc$), अतः $a \mid c$।

(2) $c = aq$ लिखिए; $b \mid aq$ और $\gcd(a, b) = 1$ से, बिंदु (1) $b \mid q$
देता है, अतः $ab \mid aq = c$।

(3) $au + bv = 1$ और $au' + cv' = 1$। दोनों संबंधों को गुणा कीजिए:
\[
1 = (au + bv)(au' + cv')
= a\,\bigl(auu' + ucv' + u'bv\bigr) + bc\,(vv') ,
\]
जो $a$ और $bc$ के बीच एक बेज़ू संबंध है: \cref{cor:b1:arith:bezout} से
$\gcd(a, bc) = 1$।
\end{proof}

\begin{example}[एक रैखिक डायोफैंटीय समीकरण हल करना]\label{ex:b1:arith:diophantine}
$6x + 10y = 4$ वाले सभी $(x, y) \in \Z^2$ खोजिए। पहले \emph{अस्तित्व की
जाँच}: $\gcd(6, 10) = 2$ $4$ को \omterm{def:b1:arith:divides}{विभाजित करता है}, अतः हल विद्यमान हैं (यदि
\omterm{thm:b1:arith:gcd}{महत्तम समापवर्तक} दाएँ पक्ष को विभाजित न करता, तो बायाँ पक्ष सदा उसका गुणज
होता और कोई हल न होता)। पूरे समीकरण को भाग दीजिए: $3x + 5y = 2$। एक विशिष्ट
हल दिखाई देता है: $(x_0, y_0) = (-1, 1)$। व्यापक हल के लिए घटाइए: $3(x + 1)
= -5(y - 1)$, अतः $3 \mid 5(y-1)$, और गाउस की प्रमेयिका ($\gcd(3,5) = 1$) $3
\mid y - 1$ देती है: $y = 1 - 3k$, फिर $x = -1 + 5k$। विलोमतः ऐसा हर युग्म
काम करता है:
\[
(x, y) = (-1 + 5k,\ 1 - 3k), \qquad k \in \Z .
\]
यह प्रतिरूप सामान्य है: एक विशिष्ट हल, और उसके साथ $\bigl(\frac b{\gcd},
-\frac a{\gcd}\bigr)$ के पूर्णांक गुणज --- वही ``विशिष्ट और समांगी'' संरचना
जो \cref{ch:b1:diffeq} में है, जहाँ अद्वितीयता की भूमिका गाउस की प्रमेयिका
निभाती है।
\end{example}

\begin{definition}[लघुत्तम समापवर्त्य]\label{def:b1:arith:lcm}
$\operatorname{lcm}(a, b)$\index{लघुत्तम समापवर्त्य} उपसमूह $a\Z \cap b\Z$
का $\N$ में जनक है: वह $a$ और $b$ का ऐसा उभयनिष्ठ गुणज है जो प्रत्येक
उभयनिष्ठ गुणज को \omterm{def:b1:arith:divides}{विभाजित करता है}, और $a, b \in \N^*$ के लिए
\[
\gcd(a,b) \times \operatorname{lcm}(a,b) = ab
\qquad (\text{उपपत्ति } \cref{exo:b1:arith:5}).
\]
\end{definition}

\begin{example}[संरेखण की समस्याएँ लघुत्तम समापवर्त्य की समस्याएँ हैं]\label{ex:b1:arith:gears}
दो जुड़े हुए दंतचक्रों के $84$ और $36$ दाँते हैं। कितने दाँतों की उभयनिष्ठ
गति के बाद वे साथ-साथ अपनी आरंभिक स्थिति में लौटते हैं? विन्यास तब दोहराता
है जब बीते हुए दाँतों की संख्या $84$ और $36$ का उभयनिष्ठ गुणज हो; पहली बार
यह
\[
\operatorname{lcm}(84, 36) = \frac{84 \times 36}{\gcd(84, 36)}
= \frac{3024}{12} = 252
\]
दाँतों पर होता है --- अर्थात् बड़े दंतचक्र के $3$ और छोटे के $7$ चक्कर
($252/84$ और $252/36$)। व्यावहारिक मार्ग पर ध्यान दीजिए: \emph{पहले \omterm{thm:b1:arith:gcd}{महत्तम
समापवर्तक} संगणित कीजिए} (यूक्लिड: $84 = 2\times36 + 12$, $36 = 3\times12$),
फिर भाग दीजिए --- गुणजों की सूची बनाकर लघुत्तम समापवर्त्य कभी मत बनाइए।
आवर्ती संपात का हर प्रश्न (दंतचक्र, ग्रहों का संरेखण, आवर्ती दशमलवों का
मिलना) इसी एक संगणना पर सिमट जाता है।
\end{example}

\section{अभाज्य संख्याएँ}

\begin{definition}\label{def:b1:arith:prime}
पूर्णांक $p \geq 2$ \emph{अभाज्य}\index{अभाज्य संख्या} तब कहलाता है जब उसके
एकमात्र धनात्मक भाजक $1$ और $p$ हों। अभाज्य $p$ और $a \in \Z$ के लिए: या तो
$p \mid a$, या $\gcd(p, a) = 1$। फलस्वरूप (\cref{thm:b1:arith:gauss})
\emph{यूक्लिड की प्रमेयिका} सत्य है: यदि $p \mid ab$, तो $p \mid a$ या $p
\mid b$।
\end{definition}

\begin{remark}[भाग देकर अभाज्यता की जाँच]\label{rem:b1:arith:trialdivision}
यदि $2 \leq a \leq b$ के साथ $n = ab$, तो $a^2 \leq ab = n$, अतः $a \leq
\sqrt n$: किसी भाज्य $n$ का सदा कोई \omterm{def:b1:arith:prime}{अभाज्य} भाजक $\leq \sqrt n$ होता है। अतः
यह जाँचने के लिए कि $n$ \omterm{def:b1:arith:prime}{अभाज्य} है या नहीं, $\sqrt n$ तक के \omterm{def:b1:arith:prime}{अभाज्य} आज़मा लेना
पर्याप्त है। $n = 271$ के लिए: $\sqrt{271} < 17$, और $271$ $2, 3, 5, 7, 11,
13$ में से किसी से विभाज्य नहीं है (विषम, अंकों का योग $10$, अंत में $0$ या
$5$ नहीं, $271 = 7\cdot38 + 5 = 11\cdot24 + 7 = 13\cdot20 + 11$): अतः
\omterm{def:b1:arith:prime}{अभाज्य}, और वह भी दो सौ के बदले छह भागों में। $\sqrt n$ की बाधा एक सच्ची
देहली है: सौ अंकों वाली संख्याओं के लिए उसे कुशलता से पार करने के लिए वे
आधुनिक अभाज्यता-परीक्षण चाहिए जो \cref{thm:b1:arith:fermat} से उपजे हैं।
\end{remark}

\begin{theorem}[यूक्लिड]\label{thm:b1:arith:euclidprimes}
\omterm{def:b1:arith:prime}{अभाज्य} संख्याएँ अनंत हैं।
\end{theorem}

\begin{proof}
प्रत्येक पूर्णांक $n \geq 2$ का कोई \omterm{def:b1:arith:prime}{अभाज्य} भाजक है: उसका लघुतम भाजक $\geq 2$
\omterm{def:b1:arith:prime}{अभाज्य} है (उसका उचित गुणनखंडन $n$ का उससे छोटा भाजक दे देता)। अब मान लीजिए
$p_1, \dots, p_k$ ही सारे \omterm{def:b1:arith:prime}{अभाज्य} हैं, और $N = p_1 p_2 \cdots p_k + 1 \geq 2$
रखिए। कोई \omterm{def:b1:arith:prime}{अभाज्य} $p_i$ $N$ को \omterm{def:b1:arith:divides}{विभाजित करता है}; पर $p_i$ $N - 1 = p_1\cdots
p_k$ को भी \omterm{def:b1:arith:divides}{विभाजित करता है}, अतः $p_i \mid 1$ --- असंगत।
\end{proof}

\begin{theorem}[अंकगणित की मूल प्रमेय]\label{thm:b1:arith:fta}
प्रत्येक पूर्णांक $n \geq 2$ अभाज्यों का गुणनफल है, और गुणनखंडन
\[
n = p_1^{\alpha_1} p_2^{\alpha_2} \cdots p_k^{\alpha_k}
\qquad (p_1 < p_2 < \dots < p_k \text{ अभाज्य},\ \alpha_i \in \N^*)
\]
अद्वितीय है।\index{अभाज्य गुणनखंडन}
\end{theorem}

\begin{proof}
\emph{अस्तित्व} प्रबल आगमन (\cref{thm:b1:logic:induction}) से: $n = 2$
\omterm{def:b1:arith:prime}{अभाज्य} है; $n > 2$ के लिए या तो $n$ \omterm{def:b1:arith:prime}{अभाज्य} है, या $2 \leq a, b < n$ के साथ
$n = ab$, और आगमन परिकल्पना $a$ तथा $b$ का गुणनखंडन कर देती है।

\emph{अद्वितीयता।} मान लीजिए $p_1 \cdots p_r = q_1 \cdots q_s$ (\omterm{def:b1:arith:prime}{अभाज्य}
पुनरावृत्ति सहित सूचीबद्ध, मान लीजिए $r \leq s$), और $r$ पर आगमन कीजिए। यदि
$r = 0$, तो बायाँ पक्ष $1$ है, जो $s = 0$ पर बाध्य कर देता है (अभाज्यों का
अरिक्त गुणनफल $1$ से बड़ा होता है)। $r \geq 1$ के लिए: \omterm{def:b1:arith:prime}{अभाज्य} $p_1$
$q_1(q_2\cdots q_s)$ को \omterm{def:b1:arith:divides}{विभाजित करता है}, अतः यूक्लिड की प्रमेयिका से या तो
$p_1 \mid q_1$ या $p_1 \mid q_2\cdots q_s$; दोहराने पर $p_1$ किसी $q_j$ को
\omterm{def:b1:arith:divides}{विभाजित करता है}। पर $q_j$ \omterm{def:b1:arith:prime}{अभाज्य} है और $p_1 \geq 2$: अनिवार्यतः $p_1 = q_j$।
इस उभयनिष्ठ गुणनखंड को काट दीजिए (यह वैध है: $\Z$ पूर्णांकीय प्रांत है),
जिससे मिलता है
\[
p_2 \cdots p_r = q_1 \cdots \widehat{q_j} \cdots q_s
\]
(टोपी का चिह्न लोप दर्शाता है), जो छोटे गुणनफलों की समता है; आगमन परिकल्पना
कहती है कि दोनों सूचियाँ $p_2, \dots, p_r$ और $q_1, \dots, \widehat{q_j},
\dots, q_s$ क्रम तक मेल खाती हैं, अतः मूल सूचियाँ भी मेल खाती थीं। घातांक
वाला रूप समान अभाज्यों को समूहबद्ध कर देता है।
\end{proof}

\begin{proposition}[मूल्यांकन]\label{prop:b1:arith:valuation}
\omterm{def:b1:arith:prime}{अभाज्य} $p$ और $n \in \N^*$ के लिए $n$ के गुणनखंडन में $p$ के घातांक को
$v_p(n)$\index{आदिक मूल्यांकन@$p$-आदिक मूल्यांकन} लिखिए (जहाँ $p \nmid n$
होने पर $v_p(n) = 0$)। तब
\[
v_p(mn) = v_p(m) + v_p(n),
\qquad
m \mid n \iff \forall p,\ v_p(m) \leq v_p(n),
\]
\[
v_p\bigl(\gcd(m,n)\bigr) = \min\bigl(v_p(m), v_p(n)\bigr),
\qquad
v_p\bigl(\operatorname{lcm}(m,n)\bigr) = \max\bigl(v_p(m),
v_p(n)\bigr).
\]
\end{proposition}

\begin{proof}
पहली सर्वसमिका इसलिए सत्य है कि गुणनखंडन गुणित होते हैं और $mn$ का गुणनखंडन
अद्वितीय है। यदि $m \mid n$, तो $n = mq$ लिखिए और इसे लगाइए। विलोमतः, यदि
सभी $v_p(m) \leq v_p(n)$, तो पूर्णांक $q = \prod_p p^{\,v_p(n) - v_p(m)}$
$mq = n$ को संतुष्ट करता है। \omterm{thm:b1:arith:gcd}{महत्तम समापवर्तक} का सूत्र: कसौटी से पूर्णांक $d
= \prod p^{\min}$ दोनों को \omterm{def:b1:arith:divides}{विभाजित करता है}, और प्रत्येक उभयनिष्ठ भाजक $c$ के
लिए सभी $p$ पर $v_p(c) \leq \min$, अतः $c \mid d$; लघुत्तम समापवर्त्य के लिए
$\max$ के साथ वही तर्क।
\end{proof}

\begin{example}[मूल्यांकनों से वर्ग और घन]\label{ex:b1:arith:squarecube}
पूर्णांक $n \geq 1$ पूर्ण वर्ग है यदि और केवल यदि प्रत्येक $v_p(n)$ सम हो
(यदि $n = m^2$, तो $v_p(n) = 2v_p(m)$; विलोमतः प्रत्येक घातांक को आधा
कीजिए)। घनों के लिए भी वही, $3$ के गुणजों के साथ। अतः $21168 = 2^4 \times
3^3 \times 7^2$ न वर्ग है ($v_3 = 3$ विषम है) और न घन ($v_2 = 4$); ऐसा लघुतम
धनात्मक पूर्णांक $m$ कि $21168\,m$ घन \emph{हो}, प्रत्येक घातांक को $3$ के
अगले गुणज तक भरकर मिलता है:
\[
m = 2^{6-4} \times 3^{3-3} \times 7^{3-2} = 2^2 \times 7 = 28,
\qquad
21168 \times 28 = 2^6\,3^3\,7^3 = (2^2 \times 3 \times 7)^3
= 84^3 .
\]
सार: गुणात्मक प्रश्न (वर्ग, घन, भाजक, \omterm{thm:b1:arith:gcd}{महत्तम समापवर्तक}, लघुत्तम समापवर्त्य)
घातांक सदिशों $(v_2, v_3, v_5, \dots)$ पर \emph{निर्देशांक-दर-निर्देशांक}
प्रश्न बन जाते हैं --- और अद्वितीय गुणनखंडन ठीक यही \omterm{def:b1:logic:statement}{कथन} है कि ये निर्देशांक
विद्यमान और सुपरिभाषित हैं।
\end{example}

\section{सर्वांगसमताएँ}

\begin{definition}\label{def:b1:arith:congruence}
$n \in \N^*$ के लिए: $a \equiv b \pmod n$\index{सर्वांगसमता} जब $n \mid a -
b$। यह ऐसा \omterm{def:b1:logic:equiv}{तुल्यता संबंध} है जो योग और गुणन के अनुकूल है: यदि $a \equiv b$ और
$a' \equiv b'$ ($n$ के सापेक्ष), तो $a + a' \equiv b + b'$, $aa' \equiv
bb'$, और $k \in \N$ के लिए $a^k \equiv b^k$।
\end{definition}

\begin{example}[नवशेष परीक्षण]\label{ex:b1:arith:castingnines}
$+$ और $\times$ के साथ अनुकूलता उतनी ही पुरानी जाँच-युक्ति है जितना व्यापार।
चूँकि $10 \equiv 1 \pmod 9$, प्रत्येक पूर्णांक $9$ के सापेक्ष अपने अंकों के
योग के सर्वांगसम है (उपपत्ति \cref{exo:b1:arith:2} में)। दावा $1234 \times
567 = 699\,678$ जाँचने के लिए: अंकों के योग $1234 \equiv 1$ और $567 \equiv
18 \equiv 0 \pmod 9$ देते हैं, अतः गुणनफल $\equiv 1 \times 0 = 0$ होना
चाहिए; और सचमुच $6 + 9 + 9 + 6 + 7 + 8 = 45 \equiv 0$। जाँच पास हो जाती है
(और गुणनफल वस्तुतः सही है)। यदि किसी ने $699\,478$ बताया होता, तो अंकों का
योग $43 \equiv 7 \not\equiv 0$ उसे तुरंत पकड़ लेता। यह परीक्षण एकपक्षीय है
--- वह भूल तभी पकड़ता है जब भूल स्वयं $9$ का गुणज न हो --- और यह छोटे पैमाने
पर ठीक \cref{ex:b1:arith:pseudoprime} वाला छद्म-अभाज्य पाठ है: \omterm{def:b1:arith:congruence}{सर्वांगसमता}
की जाँच खंडन करती है, प्रमाणित नहीं करती।
\end{example}

\begin{proposition}[$n$ के सापेक्ष व्युत्क्रमणीयता]\label{prop:b1:arith:invmod}
$a$ \emph{$n$ के सापेक्ष व्युत्क्रमणीय} है (अर्थात् किसी $b$ के लिए $ab
\equiv 1 \pmod n$) यदि और केवल यदि $\gcd(a, n) = 1$। तब प्रतिलोम $n$ के
सापेक्ष अद्वितीय होता है और विस्तारित \omterm{met:b1:arith:euclid}{यूक्लिडीय कलनविधि} से संगणित होता है।
\end{proposition}

\begin{proof}
$ab \equiv 1 \pmod n$ का अर्थ है कि किसी $k$ के लिए $ab + nk = 1$: यह एक
बेज़ू संबंध है, जो तभी विद्यमान है जब $\gcd(a,n) = 1$
(\cref{cor:b1:arith:bezout})। अद्वितीयता: यदि $ab \equiv ab' \equiv 1$, तो
$b \equiv b(ab') = (ab)b' \equiv b' \pmod n$।
\end{proof}

\begin{example}[$26$ के सापेक्ष $7$ का प्रतिलोम]\label{ex:b1:arith:inv7mod26}
चूँकि $\gcd(7, 26) = 1$, $7$ का वर्ग $26$ के सापेक्ष व्युत्क्रमणीय है।
विस्तारित यूक्लिड:
\[
26 = 3 \times 7 + 5, \qquad
7 = 1 \times 5 + 2, \qquad
5 = 2 \times 2 + 1 ,
\]
फिर उलटी दिशा में:
\[
1 = 5 - 2 \times 2 = 5 - 2(7 - 5) = 3 \times 5 - 2 \times 7
= 3(26 - 3 \times 7) - 2 \times 7 = 3 \times 26 - 11 \times 7 .
\]
अतः $7 \times (-11) \equiv 1 \pmod{26}$, अर्थात् $7^{-1} \equiv -11 \equiv
15 \pmod{26}$; जाँच: $7 \times 15 = 105 = 4 \times 26 + 1$। प्रतिलोम हाथ में
होने पर कोई भी \omterm{def:b1:arith:congruence}{सर्वांगसमता} $7x \equiv c \pmod{26}$ एक ही गुणा में हल हो जाती
है: $x \equiv 15c$। यही यांत्रिक प्रतिलोमन मापांकी अंकगणित का परिश्रमी घोड़ा
है --- और \cref{rem:b1:arith:whereused} में उल्लिखित सार्वजनिक-कुंजी
प्रोटोकॉलों का भी, जहाँ \omterm{def:b1:complex:field}{मापांक} सैकड़ों अंकों के होते हैं पर कलनविधि ठीक यही
है।
\end{example}

\begin{example}[जब गुणांक व्युत्क्रमणीय न हो]\label{ex:b1:arith:noninvertible}
$12x \equiv 8 \pmod{20}$ हल कीजिए। यहाँ $\gcd(12, 20) = 4$, अतः $12$ $20$ के
सापेक्ष व्युत्क्रमणीय नहीं है --- फिर भी समीकरण साधा जा सकता है। \omterm{def:b1:arith:congruence}{सर्वांगसमता}
कहती है $20 \mid 12x - 8$; पूरे संबंध को $4$ (तीनों सामग्रियों का भाजक) से
भाग देने पर वह $5 \mid 3x - 2$ के तुल्य है, अर्थात्
\[
3x \equiv 2 \pmod 5 .
\]
अब $\gcd(3, 5) = 1$ और $3^{-1} \equiv 2 \pmod 5$ ($3 \times 2 = 6 \equiv
1$), अतः $x \equiv 4 \pmod 5$: हल $x \equiv 4, 9, 14, 19 \pmod{20}$ हैं ---
अर्थात् $20$ के सापेक्ष \emph{चार} वर्ग, जो \omterm{thm:b1:arith:gcd}{महत्तम समापवर्तक} से मेल खाते
हैं। (यदि दायाँ पक्ष $4$ से विभाज्य न होता, मान लीजिए $12x \equiv 6
\pmod{20}$, तो कोई हल होता ही नहीं: बायाँ पक्ष सदा $\equiv 0 \pmod 4$ रहता
है।) व्यापक आकार: $ax \equiv b \pmod n$ तभी हल हो सकता है जब $\gcd(a, n)
\mid b$, और तब उसके ठीक $\gcd(a, n)$ हल-वर्ग होते हैं --- सब कुछ \omterm{thm:b1:arith:gcd}{महत्तम
समापवर्तक} से भाग दीजिए और प्रतिलोम लीजिए।
\end{example}

\begin{theorem}[फर्मा की लघु प्रमेय]\label{thm:b1:arith:fermat}
मान लीजिए $p$ \omterm{def:b1:arith:prime}{अभाज्य} है। प्रत्येक $a \in \Z$ के लिए:
\[
a^p \equiv a \pmod p,
\]
और यदि $p \nmid a$, तो $a^{p-1} \equiv 1 \pmod p$।\index{फर्मा की लघु
प्रमेय}
\end{theorem}

\begin{proof}
पहले, $1 \leq k \leq p - 1$ के लिए \omterm{def:b1:counting:objects}{द्विपद गुणांक} $\binom pk =
\frac{p!}{k!(p-k)!}$ $p$ से विभाज्य है: वस्तुतः $k!\,(p-k)!\, \binom pk =
p!$ और $p$ $p!$ को \omterm{def:b1:arith:divides}{विभाजित करता है} पर $k!(p-k)!$ से \omterm{cor:b1:arith:bezout}{सहअभाज्य} है (सभी गुणनखंड
$< p$ हैं), अतः गाउस की प्रमेयिका $p \mid \binom pk$ देती है।

अब $a \in \N$ के लिए $a^p \equiv a$ आगमन से सिद्ध कीजिए। $a = 0$ के लिए
सत्य। यदि $a^p \equiv a$, तो द्विपद प्रमेय से
\[
(a+1)^p = \sum_{k=0}^{p} \binom pk a^k
\equiv a^p + 1 \equiv a + 1 \pmod p,
\]
जहाँ बीच के सभी पद $p$ के सापेक्ष लुप्त हो जाते हैं। $a < 0$ के लिए परिणाम
को $-a$ पर लगाइए और $p = 2$ (जहाँ $x \equiv -x$) को विषम $p$ (जहाँ $(-a)^p =
-a^p$) से अलग कीजिए। अंत में, यदि $p \nmid a$, तो $a^p \equiv a$ को $p$ के
सापेक्ष $a$ के प्रतिलोम से गुणा कीजिए (\cref{prop:b1:arith:invmod})।
\end{proof}

\begin{example}[फर्मा का विलोम विफल है: $341$]\label{ex:b1:arith:pseudoprime}
फर्मा की लघु प्रमेय एक सस्ता \emph{भाज्यता} परीक्षण देती है: यदि $n$ से
\omterm{cor:b1:arith:bezout}{सहअभाज्य} किसी $a$ के लिए $a^{n-1} \not\equiv 1 \pmod n$, तो $n$ \omterm{def:b1:arith:prime}{अभाज्य} नहीं
है। क्या यह परीक्षण अभाज्यता प्रमाणित भी कर सकता है? नहीं: $n = 341 = 11
\times 31$, जो भाज्य है, और $a = 2$ लीजिए। चूँकि $2^{10} = 1024 = 3 \times
341 + 1$,
\[
2^{10} \equiv 1 \pmod{341}
\qquad\Longrightarrow\qquad
2^{340} = \bigl(2^{10}\bigr)^{34} \equiv 1 \pmod{341} :
\]
भाज्य $341$ आधार $2$ के लिए फर्मा-परीक्षण पास कर जाता है (वह लघुतम ऐसा
\emph{छद्म-अभाज्य} है)। आधार $3$ उसका मुखौटा उतार देता है ($3^{340}
\not\equiv 1$), और इसीलिए व्यावहारिक अभाज्यता-परीक्षण अनेक आधारों पर परीक्षण
चलाते हैं, कुछ परिष्करणों के साथ --- इसी विचार के औद्योगिक रूप ही
\cref{rem:b1:arith:whereused} के बड़े अभाज्यों को प्रमाणित करते हैं। उपदेश:
निहितार्थ और उसका विलोम अलग-अलग जीवन जीते हैं
(\cref{rem:b1:logic:pitfalls}), प्रमेयों के लिए भी।
\end{example}

\begin{example}[व्यावहारिक सर्वांगसमता संगणनाएँ]\label{ex:b1:arith:congruences}
$11$ के सापेक्ष $7^{2026}$ का शेषफल क्या है? फर्मा से $7^{10} \equiv 1
\pmod{11}$। चूँकि $2026 = 10 \times 202 + 6$:
\[
7^{2026} \equiv 7^6 = (7^2)^3 = 49^3 \equiv 5^3 = 125 \equiv 4
\pmod{11}.
\]
शेषफल $4$ है। रणनीति: फर्मा से मिली कोटि के सापेक्ष घातांक घटाइए, फिर हर पद
पर बीच की घातें भी घटाते जाइए।
\end{example}

\begin{remark}[अंकगणित की सामान्य भूलें]\label{rem:b1:arith:pitfalls}
\begin{enumerate}
  \item \emph{\omterm{def:b1:arith:congruence}{सर्वांगसमता} में भाग देना।} $ac \equiv bc \pmod n$ से $a \equiv
  b$ तब तक \emph{नहीं} निकाला जा सकता जब तक $\gcd(c, n) = 1$ न हो: $6 \equiv
  2 \pmod 4$, पर $3 \not\equiv 1 \pmod 4$। सही व्यापक नियम \omterm{def:b1:complex:field}{मापांक} को भी भाग
  देता है: $ac \equiv bc \pmod n \iff a \equiv b \pmod{n/\gcd(c,n)}$।
  \item \emph{यूक्लिड की प्रमेयिका का दुरुपयोग।} $a \mid bc$ से $a \mid b$
  या $a \mid c$ केवल तभी निकलता है जब $a$ \emph{\omterm{def:b1:arith:prime}{अभाज्य}} हो (या किसी एक
  गुणनखंड से \omterm{cor:b1:arith:bezout}{सहअभाज्य}): $6 \mid 4 \times 9$, फिर भी $6$ किसी भी गुणनखंड को
  विभाजित नहीं करता।
  \item \emph{\omterm{cor:b1:arith:bezout}{सहअभाज्य} होना संबंध है, गुणधर्म नहीं।} ``$8$ और $9$ \omterm{cor:b1:arith:bezout}{सहअभाज्य}
  हैं'' सत्य है, यद्यपि इनमें से कोई \omterm{def:b1:arith:prime}{अभाज्य} नहीं है; ``जोड़ों में \omterm{cor:b1:arith:bezout}{सहअभाज्य}''
  ``समग्र रूप से \omterm{cor:b1:arith:bezout}{सहअभाज्य}'' से प्रबल है ($\gcd(6, 10, 15) = 1$, पर कोई युग्म
  \omterm{cor:b1:arith:bezout}{सहअभाज्य} नहीं)।
  \item \emph{घातांक $n$ के सापेक्ष नहीं रहते।} $a^k \bmod n$ में घातांक
  केवल $a$ की \emph{कोटि} के सापेक्ष घटाया जा सकता है (उदाहरणार्थ $p - 1$,
  जब फर्मा लागू हो), कभी $n$ के सापेक्ष नहीं: $2^{10} \bmod 11$ $1$ है,
  $2^{10 \bmod 11} = 2^{10}$ नहीं --- जो घटाव काम करता है वह
  \cref{ex:b1:arith:congruences} वाला है।
\end{enumerate}
\end{remark}

\begin{remark}[यह अध्याय कहाँ काम आता है]\label{rem:b1:arith:whereused}
यह अध्याय जितना औज़ार-संदूक है उतना ही एक आदर्श भी। पूरी शृंखला ---
यूक्लिडीय भाग, \omterm{thm:b1:arith:gcd}{महत्तम समापवर्तक}, बेज़ू, गाउस, अद्वितीय गुणनखंडन ---
\cref{ch:b1:poly} में बहुपदों के लिए अक्षरशः दोहराई जाती है, जहाँ निरपेक्ष
मान की भूमिका ``घात'' निभाती है; दोनों अध्यायों को साथ-साथ रखकर देखना दोनों
को समझने का सर्वोत्तम उपाय है। सर्वांगसमताओं का कलन \cref{ch:b1:structures}
में वलय $\Z/n\Z$ बन जाता है, जिसके व्युत्क्रमणीय अवयव
(\cref{prop:b1:arith:invmod}) इकाई-समूह का पहला अतुच्छ उदाहरण बनाते हैं।
मूल्यांकन नीचे दी गई सप्ताहांत समस्या में (लजांद्र का सूत्र) लौटते हैं और
\cref{ch:b1:reals} की अपरिमेयता की उपपत्तियों को शक्ति देते हैं। इस खंड से
आगे, $n$ के सापेक्ष बेज़ू प्रतिलोमन सार्वजनिक-कुंजी गूढ़लेखन का इंजन है, और
फर्मा की लघु प्रमेय उन अभाज्यता परीक्षणों की पितामही है जो वहाँ प्रयुक्त
बड़े अभाज्यों को प्रमाणित करते हैं।
\end{remark}

\begin{remark}[मध्यांतर: $\Z$ एक आदर्श के रूप में]\label{rem:b1:arith:interlude}
अलग-अलग प्रमेयों से पीछे हटकर इस अध्याय की वास्तुकला देखिए: एक औज़ार
(यूक्लिडीय भाग) ने एक वर्गीकरण ($n\Z$ उपसमूह) दिया, उसने एक अस्तित्व प्रमेय
(\omterm{thm:b1:arith:gcd}{महत्तम समापवर्तक}, बेज़ू) दी, उसने एक विभाज्यता-कलन (गाउस) दिया, और उसने
अद्वितीय गुणनखंडन --- हर मंज़िल केवल अपने नीचे वाली पर टिकी हुई। यही भवन इस
खंड में दो बार और, भिन्न भूतलों पर, खड़ा किया जाएगा: \cref{ch:b1:poly} में,
जहाँ आकार से भाग की जगह घात से भाग आ जाता है और ऊपर का सब कुछ \emph{अक्षरशः}
दोहराता है; और लघु रूप में \cref{ch:b1:structures} के हर $\Z/n\Z$ के भीतर,
जहाँ व्युत्क्रमणीयता के प्रश्न (इस अध्याय के \cref{prop:b1:arith:invmod})
वलयों और क्षेत्रों के विषय में संरचनात्मक \omterm{def:b1:logic:statement}{कथन} बन जाते हैं। किसी तर्क को
``$\Z$-वाला तर्क, प्रतिरोपित'' के रूप में पहचान लेना उन अध्यायों को सीखने का
सबसे तेज़ उपाय है --- और बीजगणित की मूल आदत का पहला स्वाद, अर्थात् वस्तुओं
के बदले \emph{अभिगृहीतों} के विषय में प्रमेय सिद्ध करना।
\end{remark}

\begin{omfigure}
\begin{tikzpicture}[scale=0.42]
  \foreach \n/\k in {0/0, 1/0, 1/1, 2/0, 2/2, 3/0, 3/1, 3/2, 3/3,
      4/0, 4/4, 5/0, 5/1, 5/4, 5/5, 6/0, 6/2, 6/4, 6/6,
      7/0, 7/1, 7/2, 7/3, 7/4, 7/5, 7/6, 7/7}
    \fill[omDef] ({2*\k - \n}, {-\n}) circle (0.32);
  \foreach \n/\k in {2/1, 4/1, 4/2, 4/3, 5/2, 5/3, 6/1, 6/3, 6/5}
    \draw[gray] ({2*\k - \n}, {-\n}) circle (0.32);
\end{tikzpicture}

{\small पास्कल त्रिभुज की $0$ से $7$ तक की पंक्तियाँ, जिनमें \emph{विषम}
प्रविष्टियाँ भरी हुई हैं: पंक्ति $n$ में उनकी संख्या $2^{s_2(n)}$ है, जहाँ
$s_2(n)$ $n$ के द्विआधारी लेखन में इकाइयों की संख्या है (पंक्तियाँ $1, 2,
4$: दो विषम प्रविष्टियाँ; पंक्ति $7 = (111)_2$: सभी आठ)। यह स्वसमरूप
प्रतिरूप --- हर ``विषमों का त्रिभुज'' अपनी दो प्रतियाँ जनता है --- चित्र रूप
में कुमर की प्रमेय है, जो नीचे दी गई सप्ताहांत समस्या में सिद्ध की गई है।}
\end{omfigure}

\section{अभ्यास}

\begin{exercise}[$\star$]\label{exo:b1:arith:1}
\omterm{met:b1:arith:euclid}{यूक्लिडीय कलनविधि} से $\gcd(1\,001, 777)$ संगणित कीजिए, और उसके लिए एक बेज़ू
युग्म भी।
\end{exercise}

\begin{exercise}[$\star$]\label{exo:b1:arith:2}
आधार $10$ में \omterm{def:b1:arith:divides}{विभाज्यता} के नियम सिद्ध कीजिए: कोई पूर्णांक $9$ के सापेक्ष
अपने अंकों के योग के सर्वांगसम है, और $11$ के सापेक्ष अपने अंकों के एकांतरित
योग के। $9$ के सापेक्ष और $11$ के सापेक्ष $123\,456\,789$ क्या है?
\end{exercise}

\begin{exercise}[$\star$]\label{exo:b1:arith:3}
$\Z$ में हल कीजिए: $91x \equiv 1 \pmod{237}$ \emph{(विस्तारित यूक्लिड)}।
\end{exercise}

\begin{exercise}[$\star$]\label{exo:b1:arith:4}
$17x + 39y = 1$ वाले सभी युग्म $(x, y) \in \Z^2$ खोजिए; फिर $17 x + 39 y =
5$ वाले सभी युग्म।
\end{exercise}

\begin{exercise}[$\star\star$]\label{exo:b1:arith:5}
सिद्ध कीजिए कि $a, b \in \N^*$ के लिए $\gcd(a,b) \times
\operatorname{lcm}(a,b) = ab$। \emph{(\cref{prop:b1:arith:valuation} और
$\min(\alpha,\beta) + \max(\alpha,\beta) = \alpha + \beta$ के
मूल्यांकन-सूत्रों का प्रयोग कीजिए।)}
\end{exercise}

\begin{exercise}[$\star\star$]\label{exo:b1:arith:6}
मान लीजिए $a = 2^{10} \times 3^4 \times 5^2$ और $b = 2^6 \times 3^7 \times
7$। $\gcd(a, b)$, $\operatorname{lcm}(a,b)$, तथा $a$ के धनात्मक भाजकों की
संख्या संगणित कीजिए। \emph{(भाजक-गणना सूत्र $\prod_i (\alpha_i + 1)$ सिद्ध
कीजिए।)}
\end{exercise}

\begin{exercise}[$\star\star$]\label{exo:b1:arith:7}
मूल्यांकनों का प्रयोग करके सिद्ध कीजिए कि प्रत्येक \omterm{def:b1:arith:prime}{अभाज्य} $p$ के लिए $\sqrt
p$ अपरिमेय है: $p q^2 = r^2$ के दोनों पक्षों के $v_p$ की तुलना कीजिए।
\end{exercise}

\begin{exercise}[$\star\star$]\label{exo:b1:arith:8}
(चीनी शेषफल समस्या) ऐसे सभी पूर्णांक $x$ खोजिए कि
\[
x \equiv 2 \pmod 7, \qquad x \equiv 5 \pmod{11}.
\]
साथ ही सिद्ध कीजिए कि \omterm{cor:b1:arith:bezout}{सहअभाज्य} $m, n$ के लिए सर्वांगसमताओं के युग्म $x
\equiv a \ (m)$, $x \equiv b\ (n)$ का सदा हल होता है, जो $mn$ के सापेक्ष
अद्वितीय है।\index{चीनी शेषफल प्रमेय}
\end{exercise}

\begin{exercise}[$\star\star$]\label{exo:b1:arith:9}
$7$ के सापेक्ष $3^{1000}$ संगणित कीजिए, और $7^{100}$ के अंतिम दो दशमलव अंक
\emph{($100 = 4 \times 25$ के सापेक्ष: \cref{exo:b1:arith:8} का प्रयोग
कीजिए)}।
\end{exercise}

\begin{exercise}[$\star\star\star$]\label{exo:b1:arith:10}
$m, n \in \N^*$ के लिए सिद्ध कीजिए कि $\gcd(2^m - 1,\, 2^n - 1) =
2^{\gcd(m,n)} - 1$। \emph{संकेत: पहले दिखाइए कि $2^n - 1$ के सापेक्ष $2^m -
1$ का शेषफल $2^r - 1$ है, जहाँ $r$ $n$ के सापेक्ष $m$ का शेषफल है; फिर
\omterm{met:b1:arith:euclid}{यूक्लिडीय कलनविधि} का अनुसरण कीजिए।}
\end{exercise}

\begin{exercise}[$\star\star\star$]\label{exo:b1:arith:11}
(विल्सन की प्रमेय) मान लीजिए $p$ \omterm{def:b1:arith:prime}{अभाज्य} है। सिद्ध कीजिए कि
\[
(p-1)! \equiv -1 \pmod p ,
\]
--- इसके लिए $(p-1)!$ के प्रत्येक गुणनखंड को $p$ के सापेक्ष उसके प्रतिलोम के
साथ युग्मित कीजिए और स्वयं से युग्मित होने वाले गुणनखंडों की पहचान कीजिए
(पहले $x^2 \equiv 1 \pmod p$ हल कीजिए)। विलोम भी जाँचिए: यदि $n \geq 2$
\omterm{def:b1:arith:prime}{अभाज्य} नहीं है, तो $(n-1)! \not\equiv -1 \pmod n$।
\end{exercise}

\begin{exercise}[$\star\star\star$]\label{exo:b1:arith:12}
(फर्मा संख्याएँ) $n \in \N$ के लिए $F_n = 2^{2^n} + 1$ रखिए।
\begin{enumerate}
  \item सिद्ध कीजिए कि $n \geq 1$ के लिए $F_0 F_1 \cdots F_{n-1} = F_n - 2$
  (आगमन)।
  \item उससे निकालिए कि फर्मा संख्याएँ जोड़ों में \omterm{cor:b1:arith:bezout}{सहअभाज्य} हैं।
  \item उससे \cref{thm:b1:arith:euclidprimes} से स्वतंत्र यह दूसरी उपपत्ति
  निकालिए कि \omterm{def:b1:arith:prime}{अभाज्य} संख्याएँ अनंत हैं।
\end{enumerate}
\end{exercise}

\section{समस्या: लजांद्र का सूत्र और कुमर के हासिल}

\begin{problem}\label{pb:b1:arith:1}
$1000!$ के दशमलव लेखन के अंत में कितने शून्य आते हैं --- और, इससे गहरे, $n!$
को, अथवा किसी \omterm{def:b1:counting:objects}{द्विपद गुणांक} को, विभाजित करने वाली \omterm{def:b1:arith:prime}{अभाज्य} $p$ की ठीक-ठीक घात
क्या है? पूरे उत्तर प्रारंभिक अंकगणित के दो रत्न हैं: \emph{लजांद्र का
सूत्र} \index{लजांद्र का सूत्र} $v_p(n!) = \sum_{k\geq1} \lfloor n/p^k
\rfloor$, अपने अंकीय अवतार $v_p(n!) = \frac{n - s_p(n)}{p-1}$ के साथ, और
\emph{कुमर की प्रमेय}\index{कुमर की प्रमेय}: $v_p\binom{m+n}m$ आधार $p$ में
$m$ और $n$ जोड़ते समय के \emph{हासिलों} की गणना करता है। यह समस्या दोनों
सिद्ध करती है, उन्हें संख्यात्मक रूप से एक-दूसरे के सामने जाँचती है, और
चिरपरिचित परिणाम बटोरती है --- अंत के शून्य, पास्कल त्रिभुज की सम-विषमता, और
\omterm{def:b1:arith:prime}{अभाज्य संख्या} प्रमेय की दिशा में पहला परिबंध। आगे सर्वत्र $p$ \omterm{def:b1:arith:prime}{अभाज्य} है,
$\floor{x}$ पूर्णांक भाग है, और $s_p(n)$ आधार $p$ में लिखी $n$ के अंकों का
योग दर्शाता है।

\medskip \noindent\textbf{भाग I --- फ़र्श, मूल्यांकन और लजांद्र का सूत्र।}
\begin{enumerate}
  \item अभ्यास: $10!$ संगणित कीजिए और अंत के शून्यों की संख्या पढ़िए;
  $v_2(10!)$ और $v_5(10!)$ प्रत्येक गुणनखंड $1, 2, \dots, 10$ के गुणनखंडन से
  सीधे संगणित कीजिए।
  \item सिद्ध कीजिए कि $x \in \R$ और $n \in \N^*$ के लिए $\bigl\lfloor
  \lfloor x \rfloor / n \bigr\rfloor = \lfloor x/n \rfloor$।
  \item सिद्ध कीजिए कि सभी $a, b \in \N^*$ के लिए $v_p(a + b) \geq
  \min\bigl(v_p(a), v_p(b)\bigr)$, और जब $v_p(a) \neq v_p(b)$ हो तब समता।
  \item दिखाइए कि $\intint1n$ में $m$ के गुणजों की संख्या $\lfloor n/m
  \rfloor$ है।
  \item \emph{लजांद्र का सूत्र} सिद्ध कीजिए: प्रत्येक $n \in \N^*$ के लिए
        \[
        v_p(n!) = \sum_{k=1}^{\infty}
        \Bigl\lfloor \frac{n}{p^k} \Bigr\rfloor
        \]
        (यह परिमित योग है: $p^k > n$ होते ही पद लुप्त हो जाते हैं)।
        \emph{प्रत्येक $k$ के लिए $\intint1n$ के उन गुणनखंडों को गिनिए जो
        $p^k$ से विभाज्य हैं: वे जिस-जिस स्तर तक पहुँचते हैं, हर स्तर पर ठीक
        एक इकाई का योगदान करते हैं।}
\end{enumerate}

\medskip \noindent\textbf{भाग II --- अंकीय रूप और अंत के शून्य।}
\begin{enumerate}[resume]
  \item $v_5(1000!)$ और $v_2(1000!)$ संगणित कीजिए, और निष्कर्ष निकालिए:
  $1000!$ के अंत में कितने शून्य हैं?
  \item लजांद्र के सूत्र का अंकीय रूप सिद्ध कीजिए: $n = \sum_i a_i p^i$ को
  आधार $p$ में लिखने पर,
        \[
        v_p(n!) = \frac{n - s_p(n)}{p - 1} .
        \]
  \item $p = 2$ के लिए दो परिणाम: दिखाइए कि $2^n$ $n!$ को कभी विभाजित नहीं
  करता, और यह कि $2^{n-1}$ $n!$ को ठीक तभी \omterm{def:b1:arith:divides}{विभाजित करता है} जब $n$ $2$ की कोई
  घात हो।
  \item न्यूनता का परिबंध दीजिए: दिखाइए $\frac n{p-1} - \log_p(n) - 1 \leq
  v_p(n!) < \frac n{p-1}$, जिससे $\frac{v_p(n!)}{n} \to \frac1{p-1}$:
  दीर्घकाल में प्रति इकाई एक गुणनखंड $p$ का अनुपात $\frac1{p-1}$ जुड़ता जाता
  है।
  \item मान लीजिए $Z(n) = v_5(n!)$ $n!$ के अंत के शून्यों की संख्या है।
  दिखाइए $Z(n) - Z(n-1) = v_5(n)$, उससे निकालिए कि $Z$ मान $5$ को पूरी तरह
  छोड़ देता है ($Z(24)$ और $Z(25)$ संगणित कीजिए), और सिद्ध कीजिए कि किसी भी
  क्रमगुणित के अंत में ठीक पाँच शून्य नहीं होते।
\end{enumerate}

\medskip \noindent\textbf{भाग III --- कुमर की प्रमेय।}
\begin{enumerate}[resume]
  \item सिद्ध कीजिए कि सभी $x, y \in \R$ के लिए $\lfloor x + y \rfloor -
  \lfloor x \rfloor - \lfloor y \rfloor \in \{0, 1\}$, और लजांद्र के सूत्र
  से निकालिए कि
        \[
        v_p\binom{m+n}m
        = \sum_{k\geq1}\Bigl(
        \Bigl\lfloor\frac{m+n}{p^k}\Bigr\rfloor
        - \Bigl\lfloor\frac{m}{p^k}\Bigr\rfloor
        - \Bigl\lfloor\frac{n}{p^k}\Bigr\rfloor\Bigr),
        \]
        जो ऐसे पदों का योग है जिनमें से हर एक $0$ या $1$ है।
  \item \emph{कुमर की प्रमेय} सिद्ध कीजिए: उस योग का $k$-वाँ पद ठीक तब $1$
  के बराबर है जब आधार $p$ में $m$ और $n$ का योग स्थान $k$ में हासिल उत्पन्न
  करता हो; अतः $v_p\binom{m+n}m$ हासिलों की कुल संख्या है। \emph{($0 \leq
  m_0, n_0 < p^k$ के साथ $m = p^km_1 + m_0$ और $n = p^kn_1 + n_0$ लिखिए और
  $\lfloor (m_0 + n_0)/p^k \rfloor$ का निरीक्षण कीजिए।)}
  \item उससे $0 < j < p^k$ के लिए निकालिए:
        \[
        v_p\binom{p^k}{j} = k - v_p(j) ,
        \]
        जोड़ $j + (p^k - j)$ में हासिल गिनकर। (विशेष रूप से $0 < j < p$ के
        लिए $p \mid \binom p j$: \cref{thm:b1:arith:fermat} का मुख्य पद, फिर
        से प्राप्त।)
  \item सिद्ध कीजिए कि $v_2\binom{2n}n = s_2(n)$। उससे निकालिए कि केंद्रीय
  \omterm{def:b1:counting:objects}{द्विपद गुणांक} सदा सम है, और यह कि $\binom{2n}n \equiv 2 \pmod 4$ ठीक तब जब
  $n$ $2$ की कोई घात हो।
  \item वांडरमोंड सर्वसमिका (\cref{exo:b1:counting:7}) और प्रश्न 13 का
  प्रयोग करके दिखाइए कि प्रत्येक \omterm{def:b1:arith:prime}{अभाज्य} $p$ के लिए $\binom{2p}p \equiv 2
  \pmod p$।
  \item $v_3\binom{1000}{500}$ दो बार संगणित कीजिए: एक बार कुमर से (आधार $3$
  में $500$ लिखिए और $500 + 500$ में हासिल गिनिए), और एक बार लजांद्र के
  अंकीय रूप से ($s_3(500)$ और $s_3(1000)$ संगणित कीजिए); जाँचिए कि दोनों वही
  मान देते हैं।
\end{enumerate}

\medskip \noindent\textbf{भाग IV --- पास्कल त्रिभुज की सम-विषमता, और
अभाज्य-घनत्व का एक परिबंध।}
\begin{enumerate}[resume]
  \item अंक-कसौटी सिद्ध कीजिए: $\binom nk$ \emph{विषम} है यदि और केवल यदि
  $k$ का प्रत्येक द्विआधारी अंक $n$ के तदनुरूप अंक से अधिक न हो। आधार $p$
  में $p \nmid \binom nk$ के लिए अनुरूप कसौटी बताइए और सिद्ध कीजिए।
  \item उससे निकालिए कि पास्कल त्रिभुज की पंक्ति $n$ में ठीक $2^{s_2(n)}$
  विषम प्रविष्टियाँ हैं; पंक्तियों $4$ और $5$ पर सत्यापित कीजिए।
  \item उससे निकालिए कि सभी भीतरी प्रविष्टियाँ $\binom nk$ ($0 < k < n$) सम
  हैं यदि और केवल यदि $n$ $2$ की कोई घात हो।
  \item सिद्ध कीजिए कि $\binom{m+n}m$ को विभाजित करने वाली प्रत्येक \omterm{def:b1:arith:prime}{अभाज्य}
  घात अधिक से अधिक $m + n$ है: यदि $p^a \mid \binom{m+n}m$, तो $p^a \leq m +
  n$। \emph{(प्रश्न 11 के योग में कितने अशून्य पद हो सकते हैं?)}
  \item उससे निकालिए कि $\binom{2n}n$ $\operatorname{lcm}(1, 2, \dots, 2n)$
  को \omterm{def:b1:arith:divides}{विभाजित करता है}, और इसे निम्न परिबंध $\binom{2n}n \geq
  \frac{4^n}{2n+1}$ (जिसे आप सिद्ध करेंगे: पंक्ति $2n$ की $2n + 1$
  प्रविष्टियों में केंद्रीय सबसे बड़ी है) के साथ मिलाकर प्राप्त कीजिए
        \[
        \operatorname{lcm}(1, \dots, 2n) \geq \frac{4^n}{2n+1} :
        \]
        अर्थात् आरंभिक पूर्णांकों के उभयनिष्ठ गुणज \emph{चरघातांकी रूप से}
        बढ़ते हैं --- अभाज्यों की प्रचुरता की पहली परिमाणात्मक झलक।
\end{enumerate}

\medskip \noindent\textbf{भाग V --- संश्लेषण।}
\begin{enumerate}[resume]
  \item ऐसा लघुतम $n$ खोजिए कि $n!$ के अंत में कम से कम $2026$ शून्य हों।
  \emph{($Z(n) \approx n/4$ का आकलन कीजिए, फिर यथार्थ सूत्र से समायोजन
  कीजिए।)}
  \item एक अंतिम प्रति-जाँच: दिखाइए कि $7$ $\binom{100}{50}$ को विभाजित
  \emph{नहीं} करता --- पहले $50$ को आधार $7$ में लिखकर और जाँचकर कि जोड़ $50
  + 50$ में कोई हासिल नहीं है, फिर लजांद्र के सूत्र से $v_7(100!)$ और
  $v_7(50!)$ संगणित करके।
  \item इस समस्या में ठीक कहाँ प्रयोग हुआ: (क) अद्वितीय गुणनखंडन; (ख)
  यूक्लिडीय भाग का वियोजन $n = p^k n_1 + n_0$; (ग) \cref{ch:b1:counting} का
  कोई गणना-तर्क? प्रत्येक के लिए एक वाक्य।
  \item एक छोटे अनुच्छेद में संश्लेषण: लजांद्र का सूत्र \omterm{def:b1:arith:divides}{विभाज्यता} के प्रश्न
  को अंकों के अंकगणित में बदल देता है, और कुमर की प्रमेय उत्तर को एक ही जोड़
  के हासिलों से पढ़ लेती है --- इस अनुवाद पर, प्रश्न 16 की जाँचों पर, और
  प्रश्न 21 का परिबंध अभाज्यों के विषय में जो संकेत देता है उस पर टिप्पणी
  कीजिए (पूरा \omterm{def:b1:logic:statement}{कथन}, अर्थात् \omterm{def:b1:arith:prime}{अभाज्य संख्या} प्रमेय, इस खंड से बहुत परे है; इस
  अध्याय के औज़ार-संदूक का बहुपदीय अनुरूप \cref{ch:b1:poly} है)।
\end{enumerate}
\end{problem}
